
\section{Novikov series}
\label{sec:novikov_series}

The goal of this section is to define multivariable Novikov series and prove
basic facts about them.

\subsection{Multivariable Novikov series}

\begin{definition}[Definition~2.4: multivariable Novikov series]
  \label{def:novikov_series}
  \lean{Novikov.NovikovSeries}
  \leanok
  Let $A$ be an abelian group, and let $\Gamma \subseteq \mathbb{R}$ be an
  additive submonoid. We define the abelian group $A((t_1^\Gamma, \dotsc,
  t_n^\Gamma))$ to consist of coefficient functions $a\colon\Gamma^n\to A$,
  written as formal series
  \[
    \sum_{d\in\Gamma^n} a_d t_1^{d_1}\dotsm t_n^{d_n},
  \]
  such that, for every $s_1,\dotsc,s_n\in\mathbb{R}_{>0}$ and every
  $C\in\mathbb{R}$, the set
  \[
    \left\{d\in\Gamma^n \mathrel{}\middle|
      a_d\ne 0 \text{ and } \sum_{j=1}^n s_jd_j<C\right\}
  \]
  is finite.
\end{definition}

\begin{lemma}[Lemma~2.8: exactness of Novikov series]
  \label{lem:novikov_exact}
  \lean{Novikov.map_injective, Novikov.map_surjective, Novikov.map_exact}
  \leanok
  \uses{def:novikov_series}
  The functor $A \mapsto A((t_1^\Gamma, \dotsc, t_n^\Gamma))$ is exact. That is,
  if $0 \to A \xrightarrow{f} B \xrightarrow{g} C \to 0$ is an exact sequence of
  abelian groups, then the induced sequence
  \[ 0 \to A((\underline{t}^\Gamma)) \xrightarrow{f} B((\underline{t}^\Gamma))
  \xrightarrow{g} C((\underline{t}^\Gamma)) \to 0 \]
  is also exact.
\end{lemma}

\begin{proof}
  \leanok
  Injectivity is checked coefficientwise.  For surjectivity, let
  $c\in C((\underline{t}^\Gamma))$.  For every exponent $d$ with $c_d\ne0$,
  choose $b_d\in B$ such that $g(b_d)=c_d$, and set $b_d=0$ when $c_d=0$.
  Then $\operatorname{supp}(b)\subseteq\operatorname{supp}(c)$, so $b$ still
  satisfies the Novikov finiteness condition and maps to $c$.

  Now suppose $b\in B((\underline{t}^\Gamma))$ maps to zero.  Exactness of the
  original sequence gives $b_d\in\operatorname{im}(f)$ for every $d$.  Choose
  $a_d$ with $f(a_d)=b_d$, again choosing $a_d=0$ whenever $b_d=0$.  The
  inclusion $\operatorname{supp}(a)\subseteq\operatorname{supp}(b)$ makes $a$
  a Novikov series, and $f(a)=b$.  Thus the image equals the kernel.
\end{proof}

\begin{definition}[Convolution of Novikov series]
  \label{def:novikov_multiplication}
  \lean{Novikov.novikovSeriesMul}
  \leanok
  \uses{def:novikov_series}
  Let $A$, $B$, and $C$ be abelian groups, and let
  $\alpha\colon A\times B\to C$ be a bilinear map. For
  $f\in A((t_1^\Gamma,\dotsc,t_n^\Gamma))$ and
  $g\in B((t_1^\Gamma,\dotsc,t_n^\Gamma))$, we define the product
  $f\ast_\alpha g$ by
  bilinearly extending the product
  \[ (a t_1^{d_1} \dotsm t_n^{d_n}) \ast_\alpha (b t_1^{e_1} \dotsm t_n^{e_n})
  = \alpha(a, b) t_1^{d_1+e_1} \dotsm t_n^{d_n+e_n}. \]
\end{definition}

\begin{theorem}[Novikov series form a ring]
  \label{thm:novikov_ring}
  \lean{Novikov.NovikovSeriesRing}
  \leanok
  \uses{def:novikov_series, def:novikov_multiplication}
  If $A$ is a commutative ring, then the set of Novikov series $A((t_1^\Gamma,
  \dotsc, t_n^\Gamma))$ forms a commutative ring under addition and the
  natural product.
\end{theorem}

\begin{proof}
  \leanok
  For each exponent $d$, the Novikov condition implies that only finitely many
  pairs
  \[
    (d_1,d_2)\in\operatorname{supp}(f)\times\operatorname{supp}(g),
    \qquad d_1+d_2=d,
  \]
  occur.  Hence the convolution coefficient
  $(f*g)_d=\sum_{d_1+d_2=d}f_{d_1}g_{d_2}$ is a finite sum.

  Fix positive weights $s$ and a bound $C$.  If $(f*g)_d\ne0$ and
  $\sum_j s_jd_j<C$, then some contributing pair satisfies
  \[
    \sum_j s_j(d_{1,j}+d_{2,j})<C.
  \]
  Each support has a lower bound for its $s$-weighted degree: choose one
  supported exponent and use finiteness below its degree.  A pair as above
  must therefore place each component in a fixed finite truncation of its
  respective support.  Hence there are only finitely many such pairs, and
  consequently only finitely many possible sums $d$.  Thus the product again
  satisfies the Novikov condition.  The ring laws follow by
  comparing coefficients and reindexing finite sums over pairs and triples;
  the unit is the series supported at exponent zero with coefficient $1$.
\end{proof}

\begin{theorem}[Module structure on Novikov series]
  \label{thm:novikov_module}
  \lean{Novikov.moduleNovikovSeries}
  \leanok
  \uses{thm:novikov_ring, def:novikov_multiplication}
  If $A$ is a commutative ring and $M$ is an $A$-module, then $M((t_1^\Gamma,
  \dotsc, t_n^\Gamma))$ is naturally a module over the Novikov ring
  $A((t_1^\Gamma, \dotsc, t_n^\Gamma))$.
\end{theorem}

\begin{proof}
  \leanok
  The $A((\underline{t}^\Gamma))$-module structure on $M((\underline{t}^\Gamma))$ is defined via the
  convolution product using the scalar multiplication $\alpha \colon A \times M \to M$.
  The module axioms, such as $(f \ast g) \cdot m = f \cdot (g \cdot m)$ and distributivity,
  are verified by comparing coefficients, which involves reordering sums based on the associativity
  and distributivity of the underlying module $M$ over $A$.
\end{proof}

\begin{theorem}[Base change for Novikov series]
  \label{thm:novikov_base_change}
  \lean{Novikov.novikovModule_base_change_equiv}
  \leanok
  \uses{thm:novikov_module, lem:novikov_exact}
  Let $M$ be a finitely presented $A$-module. The natural map
  \[ M \otimes_A A((t_1^\Gamma, \dotsc, t_n^\Gamma)) \to M((t_1^\Gamma, \dotsc, t_n^\Gamma)) \]
  is an isomorphism of $A((t_1^\Gamma, \dotsc, t_n^\Gamma))$-modules.
\end{theorem}

\begin{proof}
  \leanok
  Choose a finite presentation
  \[
    A^m\longrightarrow A^n\longrightarrow M\longrightarrow0.
  \]
  Applying Novikov series gives a right exact sequence by
  Lemma~\ref{lem:novikov_exact}; tensoring with
  $A((\underline{t}^\Gamma))$ gives another right exact sequence.  The natural
  base-change maps form a commutative diagram between these sequences, and
  they are isomorphisms on the two finite free modules.  Consequently the two
  rightmost terms are the cokernels of the same map under these
  identifications.  The induced map
  \[
    M\otimes_A A((\underline{t}^\Gamma))
      \longrightarrow M((\underline{t}^\Gamma))
  \]
  is therefore an isomorphism.
\end{proof}

\begin{definition}[Substitution of variables]
  \label{def:novikov_substitute}
  \lean{Novikov.substitute}
  \leanok
  \uses{def:novikov_series}
  Given a map $f \colon [n] \to [m]$ between finite sets of variables, we
  define the substitution map
  \[
    f_\ast \colon A((t_1^\Gamma, \dotsc, t_n^\Gamma)) \longrightarrow
      A((s_1^\Gamma, \dotsc, s_m^\Gamma))
  \]
  by sending $t_i$ to $s_{f(i)}$, or equivalently by pushing exponent vectors
  forward along $f$. When it is clear from the context, we shall denote this
  induced map by $f$ instead of $f_\ast$.
\end{definition}

\begin{theorem}[Functoriality of substitution]
  \label{thm:novikov_substitute_functorial}
  \lean{Novikov.substitute_id, Novikov.substitute_comp, Novikov.substituteRingHom}
  \leanok
  \uses{def:novikov_substitute, thm:novikov_ring}
  Substitution is functorial with respect to maps of finite variable sets.
  That is, $\text{id}_*=\text{id}$ and $(g\circ f)_*=g_*\circ f_*$. Furthermore,
  if $A$ is a
  commutative ring, then $f_*$ is a ring homomorphism.
\end{theorem}

\begin{proof}
  \leanok
  For an exponent vector $d\in\Gamma^n$, let
  \[
    (f_{!}d)_j=\sum_{i:\,f(i)=j}d_i.
  \]
  The coefficient of $f_*a$ at $d'\in\Gamma^m$ is the finite sum of the
  coefficients $a_d$ over $f_{!}d=d'$.  Finiteness follows by using unit weights:
  on this fiber, $\sum_i d_i=\sum_jd'_j$, so the relevant part of the support
  lies in a fixed Novikov half-space.  More generally, positive target weights
  $s'_j$ pull back to the positive source weights $s_i=s'_{f(i)}$; this proves
  the Novikov condition for the substituted series.

  The identities $(\mathrm{id})_{!}=\mathrm{id}$ and
  $(g\circ f)_{!}=g_{!}\circ f_{!}$ prove functoriality after regrouping the
  finite coefficient sums.  Finally, $f_{!}(d+e)=f_{!}d+f_{!}e$, so the same finite-sum
  reindexing shows that substitution preserves convolution and hence defines
  a ring homomorphism.
\end{proof}

\subsection{One-variable Novikov series}

When there is a single variable $t$, we give $A((t^\Gamma))$ a topology.

\begin{definition}[The Novikov filtration]
  \label{def:novikov_filtration}
  \lean{Novikov.filtration}
  \leanok
  \uses{def:novikov_series}
  For $D \in \mathbb{R}$, we define the filtration subgroup $F_D A((t^\Gamma))$
  as the set of series $f = \sum_{d \in \Gamma} a_d t^d$ such that $a_d = 0$ for
  all $d < D$.
\end{definition}

\begin{theorem}[Definition~2.15: the Novikov topology]
  \label{thm:novikov_topology}
  \lean{Novikov.is_topological_add_group,
    Novikov.instCompleteSpaceOneVarNovikovSeries,
    Novikov.is_topological_ring}
  \leanok
  \uses{def:novikov_filtration, thm:novikov_ring}
  If $A$ is an abelian group, the collection of filtration subgroups
  $\{F_D A((t^\Gamma))\}_{D \in \mathbb{R}}$ forms a filter basis for a topology
  on $A((t^\Gamma))$. Under this topology, $A((t^\Gamma))$ is a complete
  topological abelian group. Furthermore, if $A$ is a commutative ring, then
  $A((t^\Gamma))$ is a topological ring.
\end{theorem}

\begin{proof}
  \leanok
  The filtration subgroups satisfy $F_{D_2}\subseteq F_{D_1}$ for
  $D_1\le D_2$, and each $F_D$ is an additive subgroup.  They therefore form a
  neighborhood basis at zero for a topological additive group.

  To prove completeness, let $\mathcal F$ be a Cauchy filter.  For every
  exponent $d$, the Cauchy condition modulo $F_D$ for any $D>d$ shows that the
  $d$-coefficient is eventually constant on $\mathcal F$; define $a_d$ to be
  this eventual value.  Given a degree bound, one may choose a single member
  of the filter with which the resulting series agrees below that bound.
  Its support there is therefore contained in the support of an existing
  Novikov series and is finite.  Thus $a=(a_d)$ is a Novikov series, and the
  same coefficient-stability argument shows that $\mathcal F$ converges to
  $a$.

  Finally, $F_{D_1}F_{D_2}\subseteq F_{D_1+D_2}$.  Every series belongs to
  some filtration level.  Given $x,y$ and a target neighborhood $F_D$, choose
  sufficiently high filtration levels for perturbations $u,v$ so that
  $xv$, $uy$, and $uv$ all lie in $F_D$; then
  \[
    (x+u)(y+v)-xy=xv+uy+uv\in F_D.
  \]
  This proves continuity of multiplication and makes the Novikov ring a
  topological ring.
\end{proof}

\begin{theorem}[One-variable Novikov series form a field]
  \label{thm:novikov_field}
  \lean{Novikov.novikovField}
  \leanok
  \uses{thm:novikov_topology}
  If $A$ is a field and $\Gamma\subseteq\mathbb{R}$ is an additive subgroup,
  then the ring of one-variable Novikov series $A((t^\Gamma))$ is a field.
\end{theorem}

\begin{proof}
  \leanok
  The Novikov condition gives every nonzero series $f$ a minimum degree $d$
  and a nonzero leading coefficient $a=f_d$.  Because $\Gamma$ is a subgroup,
  the monomial $t^{-d}$ is available.  Normalize $f$ by setting
  \[
    h=a^{-1}t^{-d}f,\qquad g=1-h.
  \]
  Then $h$ has constant coefficient $1$, while $g$ is supported in strictly
  positive degrees.  Its powers tend to zero in the Novikov topology, and
  completeness makes the geometric series $\sum_{n\ge0}g^n$ summable.  The
  identity
  \[
    (1-g)\sum_{n\ge0}g^n=1
  \]
  shows that
  $a^{-1}t^{-d}\sum_{n\ge0}g^n$ is an inverse of $f$.
\end{proof}

\subsection{Canonical module topology}

\begin{definition}[Canonical module topology]
  \label{def:canonical_topology}
  \lean{Novikov.Miscellany.canonicalTopology}
  \leanok
  Let $A$ be a topological commutative ring, and let $M$ be an $A$-module.
  The \emph{canonical topology} on $M$ is the initial topology with respect to
  all $A$-linear maps $M \to A$. That is, it is the coarsest topology on $M$
  for which every $A$-linear map $f \colon M \to A$ is continuous.
\end{definition}

\begin{theorem}[Basic properties of the canonical topology]
  \label{thm:canonical_topology_properties}
  \lean{Novikov.Miscellany.canonicalTopology.continuousAdd,
    Novikov.Miscellany.canonicalTopology.continuousSMul,
    Novikov.Miscellany.canonicalTopology_self_eq,
    Novikov.Miscellany.canonicalTopology.continuous_linearMap,
    Novikov.Miscellany.canonicalTopology_prod_eq}
  \leanok
  \uses{def:canonical_topology}
  Let $A$ be a topological commutative ring, and let $M$, $N$ be $A$-modules.
  \begin{enumerate}
    \item The canonical topology makes $M$ into a topological $A$-module.
    \item On $A$ itself, the canonical topology coincides with the original
      topology.
    \item Every $A$-linear map $M \to N$ is continuous with respect to the
      canonical topologies.
    \item The canonical topology on $M \times N$ coincides with the product
      topology of the canonical topologies on $M$ and $N$.
  \end{enumerate}
\end{theorem}

\begin{proof}
  \leanok
  (1) By the defining initial property, continuity may be tested after
  composing with each $f\colon M\to A$.  For addition, use
  \[
    f(x+y)=f(x)+f(y),
  \]
  together with continuity of $f$ on each factor and continuity of addition
  on $A$.  Negation is similar.  For scalar multiplication, the composite
  $A\times M\to M\xrightarrow{f}A$ is
  $(a,x)\mapsto a f(x)$, which factors through
  $\mathrm{id}_A\times f\colon A\times M\to A\times A$ followed by
  multiplication on $A$.

  (2) The identity map $\mathrm{id}_A\colon A\to A$ is $A$-linear, so the
  canonical topology on $A$ is finer than the original topology.  Conversely,
  every $A$-linear map $f\colon A\to A$ is given by $f(x)=f(1)x$, which is
  continuous in the original topology because $A$ is a topological ring.
  Hence the two topologies coincide.

  (3) For any $A$-linear $f\colon N\to A$, the composition
  $f\circ\varphi\colon M\to A$ is $A$-linear and therefore continuous in the
  canonical topology on $M$.  The initial property of the canonical topology
  on $N$ then makes $\varphi$ continuous.

  (4) Any $A$-linear map $f\colon M\times N\to A$ decomposes as
  $f(x,y)=f_1(x)+f_2(y)$, where $f_1=f\circ\iota_M$ and
  $f_2=f\circ\iota_N$.  Thus every such $f$ is continuous for the product
  topology, giving one comparison.  Conversely, the two projections are
  linear and hence continuous for the canonical topology, giving the reverse
  comparison.
\end{proof}

\begin{lemma}[Split descent of finite projectivity]
  \label{lem:split_base_change_descent}
  \lean{Novikov.Miscellany.finite_projective_of_split_baseChange}
  \leanok
  Let $A$ be a commutative ring, let $M$ be an $A$-module, and let $B$ be a
  commutative $A$-algebra whose structure map admits an $A$-linear retraction
  $B\to A$.  If $B\otimes_A M$ is a finite projective $B$-module, then $M$ is
  a finite projective $A$-module.
\end{lemma}

\begin{proof}
  \leanok
  Write $\sigma\colon B\to A$ for the $A$-linear retraction.  Choose finite
  $B$-module generators $z_1,\ldots,z_r$ of $B\otimes_A M$, and express each
  one as a finite sum
  \[
    z_j=\sum_k b_{jk}\otimes m_{jk}.
  \]
  If $m\in M$, write $1\otimes m=\sum_jc_jz_j$.  Applying
  $\sigma\otimes\mathrm{id}_M$ gives
  \[
    m=\sum_{j,k}\sigma(c_jb_{jk})m_{jk},
  \]
  so the finite family $(m_{jk})$ generates $M$ over $A$.

  For projectivity, let $\varepsilon\colon P\twoheadrightarrow Q$ be an
  $A$-linear surjection and let $u\colon M\to Q$.  Projectivity of
  $B\otimes_A M$ lifts the base change of $u$ through
  $B\otimes_A\varepsilon$.  Composing this lift with
  $m\mapsto1\otimes m$ and then with
  $\sigma\otimes\mathrm{id}_P$ gives an $A$-linear lift of $u$ through
  $\varepsilon$.  Hence $M$ is projective as well as finite.
\end{proof}

\begin{lemma}[Finite generation from real Novikov series]
  \label{lem:finite_generation_real_novikov_series}
  \lean{Novikov.module_finite_of_realNovikovSeries}
  \leanok
  \uses{thm:novikov_topology}
  Let $A$ be a commutative ring, let $M$ be an $A$-module, and put
  \[
    R=A((t^{\mathbb{R}})),\qquad \widetilde M=M((t^{\mathbb{R}})).
  \]
  Suppose that $\widetilde M$ is a finite projective $R$-module and that every
  $R$-linear functional $\widetilde M\to R$ is continuous for the Novikov
  filtration topologies.  Then $M$ is finitely generated over $A$.
\end{lemma}

\begin{proof}
  \leanok
  Choose a split surjection and a section
  \[
    \pi\colon R^n\twoheadrightarrow\widetilde M,
    \qquad
    \sigma\colon\widetilde M\longrightarrow R^n,
    \qquad \pi\sigma=\mathrm{id}.
  \]
  Each component of $\sigma$ is an $R$-linear functional and hence is
  continuous.  Since there are only finitely many components, there is a
  common $N\in\mathbb{R}$ such that
  \[
    \sigma(F_N\widetilde M)\subseteq(F_0R)^n.
  \]
  Let $x_i=\pi(e_i)$.  For each $i$, only finitely many coefficients of $x_i$
  occur in degrees at most $N$; let $G\subseteq M$ be the union of these
  finite sets of coefficients.

  Given $m\in M$, apply the splitting to the monomial
  $mt^N\in F_N\widetilde M$ and write
  \[
    mt^N=\sum_i\sigma_i(mt^N)x_i.
  \]
  Every coefficient series $\sigma_i(mt^N)$ lies in $F_0R$.  Consequently,
  the coefficient at degree $N$ on the right is an $A$-linear combination of
  coefficients of the $x_i$ in degrees at most $N$, all belonging to $G$.
  The coefficient on the left is $m$, so $G$ generates $M$.
\end{proof}

\begin{lemma}[Finite presentation from real Novikov series]
  \label{lem:finite_presentation_real_novikov_series}
  \lean{Novikov.module_finitePresentation_of_realNovikovSeries}
  \leanok
  \uses{lem:finite_generation_real_novikov_series, lem:novikov_exact,
    thm:novikov_module}
  Under the hypotheses of the preceding lemma, the $A$-module $M$ is finitely
  presented.
\end{lemma}

\begin{proof}
  \leanok
  The preceding lemma makes $M$ finitely generated.  Choose a surjection
  $p\colon A^m\twoheadrightarrow M$ and set $K=\ker(p)$.  Applying real
  Novikov series gives a surjection
  \[
    \widetilde p\colon R^m\twoheadrightarrow\widetilde M
  \]
  whose kernel is canonically $K((t^{\mathbb{R}}))$, by exactness of Novikov
  series.  Projectivity of $\widetilde M$ splits $\widetilde p$, so
  $K((t^{\mathbb{R}}))$ is a finite projective direct summand of $R^m$.

  Every $R$-linear functional on this kernel extends along the direct-summand
  projection to a functional on $R^m$.  Functionals on the finite free module
  $R^m$ are continuous, and the standard coefficientwise identification
  $(A^m)((t^{\mathbb{R}}))\cong R^m$ is a homeomorphism.  Thus the continuity
  hypothesis of
  Lemma~\ref{lem:finite_generation_real_novikov_series} holds for $K$.
  Applying that lemma shows that $K$ is finitely generated, so $M$ is finitely
  presented.
\end{proof}

\begin{theorem}[Lemma~2.16: finite projectivity descends from Novikov series]
  \label{thm:projective_novikov_series}
  \lean{Novikov.projective_of_realNovikovSeries}
  \leanok
  \uses{lem:finite_presentation_real_novikov_series,
    lem:split_base_change_descent, thm:novikov_base_change}
  Let $A$ be a commutative ring and let $M$ be an $A$-module. If
  $M((t^{\mathbb{R}}))$ is a finite projective
  $A((t^{\mathbb{R}}))$-module and its canonical topology coincides with the
  natural filtration topology, then $M$ is a finite projective $A$-module.
\end{theorem}

\begin{proof}
  \leanok
  Equality of the canonical and filtration topologies makes every
  $A((t^{\mathbb{R}}))$-linear functional continuous.  Hence
  Lemma~\ref{lem:finite_presentation_real_novikov_series} shows that $M$ is
  finitely presented.  The base-change isomorphism of
  Theorem~\ref{thm:novikov_base_change} now identifies
  \[
    A((t^{\mathbb{R}}))\otimes_A M\cong M((t^{\mathbb{R}})),
  \]
  so this scalar extension is finite projective.  Finally, extraction of the
  degree-zero coefficient gives an $A$-linear retraction
  $A((t^{\mathbb{R}}))\to A$ of the constant-series inclusion.  Apply
  Lemma~\ref{lem:split_base_change_descent} to conclude that $M$ is finite
  projective over $A$.
\end{proof}
